% ============================================================
% 01-bezier-algebra-code-explained.tex
% Compile with:  pdflatex -shell-escape 01-bezier-algebra-code-explained.tex
% Run twice for table of contents.
% ============================================================

\documentclass[12pt, a4paper]{article}

% ---- Encoding & fonts ----------------------------------------
\usepackage[T1]{fontenc}
\usepackage[utf8]{inputenc}
\usepackage{lmodern}

% ---- Page geometry -------------------------------------------
\usepackage[
  top=2.5cm, bottom=2.5cm,
  left=2.8cm, right=2.8cm
]{geometry}

% ---- Colours -------------------------------------------------
\usepackage[dvipsnames]{xcolor}
\definecolor{codebg}{RGB}{248,248,248}
\definecolor{rulecolor}{RGB}{180,180,180}
\definecolor{examplebg}{RGB}{240,248,255}
\definecolor{examplerule}{RGB}{100,149,237}
\definecolor{notebg}{RGB}{255,250,230}
\definecolor{noterule}{RGB}{210,180,50}

% ---- Minted --------------------------------------------------
\usepackage{minted}
\setminted[julia]{
  bgcolor    = codebg,
  frame      = leftline,
  framerule  = 1.5pt,
  rulecolor  = rulecolor,
  linenos    = false,
  fontsize   = \small,
  breaklines = true,
  tabsize    = 4,
  style      = friendly,
}
\setminted[text]{
  bgcolor    = codebg,
  fontsize   = \small,
  breaklines = true,
}

% ---- Math ----------------------------------------------------
\usepackage{amsmath}
\usepackage{amssymb}

% ---- Boxes ---------------------------------------------------
\usepackage{mdframed}

\newmdenv[
  backgroundcolor  = examplebg,
  linecolor        = examplerule,
  linewidth        = 1.5pt,
  leftline         = true, rightline = false,
  topline          = false, bottomline = false,
  innerleftmargin  = 10pt, innerrightmargin = 8pt,
  innertopmargin   = 6pt,  innerbottommargin = 6pt,
  skipabove = 6pt, skipbelow = 4pt,
]{examplebox}

\newmdenv[
  backgroundcolor  = notebg,
  linecolor        = noterule,
  linewidth        = 1.5pt,
  leftline         = true, rightline = false,
  topline          = false, bottomline = false,
  innerleftmargin  = 10pt, innerrightmargin = 8pt,
  innertopmargin   = 6pt,  innerbottommargin = 6pt,
  skipabove = 6pt, skipbelow = 4pt,
]{notebox}

% ---- Hyperlinks ----------------------------------------------
\usepackage[colorlinks=true, linkcolor=black, urlcolor=MidnightBlue]{hyperref}

% ---- Misc ----------------------------------------------------
\usepackage{booktabs}
\usepackage{needspace}
\usepackage{tikz}
\usetikzlibrary{shapes.geometric, arrows.meta, positioning, fit, backgrounds}
\usepackage{parskip}
\usepackage{titlesec}
\usepackage{fancyhdr}
\usepackage{datetime2}

% ---- Section formatting --------------------------------------
\titleformat{\section}
  {\large\bfseries\color{MidnightBlue}}
  {\thesection.}{0.6em}{}[\vspace{-4pt}\rule{\linewidth}{0.4pt}]
\titleformat{\subsection}
  {\normalsize\bfseries\color{BrickRed}}
  {\thesubsection.}{0.5em}{}
\titleformat{\subsubsection}
  {\normalsize\bfseries}
  {\thesubsubsection.}{0.5em}{}

% ---- Header / footer -----------------------------------------
\pagestyle{fancy}
\fancyhf{}
\rhead{\textcolor{gray}{\small 01-bezier-algebra-code-explained}}
\lhead{\textcolor{gray}{\small B\'ezier --- Code Explanation}}
\cfoot{\thepage}
\renewcommand{\headrulewidth}{0.3pt}

% ---- Helpers -------------------------------------------------
\newcommand{\cd}[1]{\texttt{#1}}
\newcommand{\jl}[1]{\mintinline{julia}{#1}}

% ---- Title ---------------------------------------------------
\title{%
  \vspace{1cm}
  {\LARGE\bfseries B\'ezier Curve --- Algebraic Expansion}\\[0.4em]
  {\large Code Explanation for Julia Beginners}\\[0.4em]
  {\normalsize\texttt{01-bezier-algebra.jl}}
}
\author{E.R.\ Nurwijayadi}
\date{\DTMtoday}

\newcommand{\docnote}{\textit{Directed and curated by E.R.\ Nurwijayadi.
Code and explanations AI-generated.}}


% =============================================================
\begin{document}
% =============================================================

\maketitle
\thispagestyle{empty}

\begin{center}\docnote\end{center}

\begin{center}
\begin{minipage}{0.85\textwidth}
\small
This document explains the Julia code in \cd{01-bezier-algebra.jl}
for readers who are new to Julia.
It focuses on the parts of the language that may look unfamiliar:
type annotations, constructors, operator overriding, multiple dispatch,
and how data flows through the print functions.
\end{minipage}
\end{center}

\vspace{1em}\hrule\vspace{1.5em}
\tableofcontents
\newpage


% =============================================================
\needspace{8\baselineskip}
\section{Type Map: How the Pieces Fit Together}
% =============================================================

\textit{Before diving into individual functions,
it helps to see the whole picture at once.
Julia has no classes, but the code has a clear object-like structure:
two structs, a cluster of free functions that operate on them,
and one entry point that ties everything together.}

\textit{The diagram below shows every type and function group,
and the arrows show who calls, owns, or produces whom.}

\bigskip

% TikZ styles
\tikzset{
  structbox/.style={
    rectangle, rounded corners=4pt,
    draw=MidnightBlue, line width=1.5pt,
    fill=examplebg,
    text width=3.0cm, align=center,
    font=\small\bfseries,
    inner sep=8pt
  },
  funcbox/.style={
    rectangle, rounded corners=2pt,
    draw=gray!60, line width=1pt,
    fill=white,
    text width=3.2cm, align=center,
    font=\small,
    inner sep=7pt
  },
  entrybox/.style={
    rectangle, rounded corners=4pt,
    draw=BrickRed, line width=1.5pt,
    fill=notebg,
    text width=2.4cm, align=center,
    font=\small\bfseries,
    inner sep=8pt
  },
  myarrow/.style={
    -Stealth, thick, draw=gray!70
  },
  bluearrow/.style={
    -Stealth, thick, draw=MidnightBlue
  },
  redarrow/.style={
    -Stealth, thick, draw=BrickRed!80
  },
  arrowlabel/.style={
    font=\footnotesize\itshape, fill=white,
    inner sep=2pt
  }
}

\begin{center}
\begin{tikzpicture}[node distance=1.4cm and 2.6cm]

  % - Structs (top row) -
  \node[structbox] (poly) {%
    \texttt{Poly}\\[3pt]
    \footnotesize\normalfont
    \texttt{coeffs::Vector\{Float64\}}
  };

  \node[structbox, right=3.8cm of poly] (bezier) {%
    \texttt{BezierExpansion}\\[3pt]
    \footnotesize\normalfont
    \texttt{basis::Vector\{Poly\}}\\
    \texttt{polys::Dict\{String,Poly\}}\\
    \texttt{points, axes, n, degree, dim}
  };

  % - expand_bernstein (below Poly) -
  \node[funcbox, below=1.6cm of poly] (expand) {%
    \texttt{expand\_bernstein}\\[2pt]
    \footnotesize\normalfont
    \textit{produces a \texttt{Poly}}\\
    \textit{from} $(n, i)$
  };

  % - Poly operations (below expand) -
  \node[funcbox, below=1.2cm of expand] (polyops) {%
    \textbf{Poly operations}\\[2pt]
    \footnotesize\normalfont
    \texttt{*}\quad\texttt{+}\quad\texttt{(p)(u)}\quad\texttt{degree}
  };

  % - BezierExpansion constructor (below bezier) -
  \node[funcbox, below=1.6cm of bezier] (bezierctr) {%
    \texttt{BezierExpansion()}\\
    \texttt{evaluate()}\\[2pt]
    \footnotesize\normalfont
    \textit{constructor + evaluator}
  };

  % - Formatting helpers (bottom left) -
  \node[funcbox, below=1.2cm of polyops] (fmt) {%
    \textbf{Formatting helpers}\\[2pt]
    \footnotesize\normalfont
    \texttt{fmt\_poly}\\
    \texttt{fmt\_coeff}\\
    \texttt{fmt\_bernstein\_symbolic}
  };

  % - Print functions (bottom right) -
  \node[funcbox, below=1.2cm of bezierctr] (print) {%
    \textbf{Print functions}\\[2pt]
    \footnotesize\normalfont
    \texttt{print\_step1} \ldots \texttt{print\_step5}\\
    \texttt{print\_all}
  };

  % - main (far right) -
  \node[entrybox, right=2.0cm of bezier] (main) {%
    \texttt{main()}\\[2pt]
    \footnotesize\normalfont\normalcolor
    \textit{entry point}
  };

  % - Arrows -

  % expand_bernstein produces Poly
  \draw[bluearrow] (expand) -- node[arrowlabel, right]{produces}  (poly);

  % BezierExpansion owns Poly (via fields)
  \draw[bluearrow] (bezier) -- node[arrowlabel, above]{owns}  (poly);

  % BezierExpansion constructor calls expand_bernstein
  \draw[myarrow] (bezierctr) -- node[arrowlabel, right]{calls} (expand);

  % BezierExpansion constructor calls Poly ops
  \draw[myarrow] (bezierctr) -- node[arrowlabel, left, pos=0.4]{uses} (polyops);

  % Poly ops act on Poly
  \draw[myarrow] (polyops) -- node[arrowlabel, right]{acts on} (poly);

  % BezierExpansion constructor builds BezierExpansion
  \draw[bluearrow] (bezierctr) -- node[arrowlabel, right]{builds} (bezier);

  % fmt helpers read Poly
  \draw[myarrow] (fmt) -- node[arrowlabel, right]{reads} (polyops);

  % print functions read BezierExpansion
  \draw[myarrow] (print) -- node[arrowlabel, right]{reads} (bezierctr);

  % print functions call fmt helpers
  \draw[myarrow] (print) -- node[arrowlabel, above]{calls} (fmt);

  % main calls BezierExpansion constructor
  \draw[redarrow] (main) -- node[arrowlabel, above]{calls} (bezier);

  % main calls print_all
  \draw[redarrow] (main) |- node[arrowlabel, right, pos=0.8]{calls} (print);

\end{tikzpicture}
\end{center}

\bigskip

\begin{notebox}
\textbf{Reading the diagram:}

\medskip
\begin{tabular}{lp{10cm}}
\toprule
Arrow colour & Meaning \\
\midrule
Blue  & Structural relationship --- one type owns or produces another \\
Grey  & Function call or data dependency \\
Red   & Entry point --- \texttt{main()} triggering the whole chain \\
\bottomrule
\end{tabular}

\medskip
The two \textbf{blue boxes} are structs (Julia types).
Everything else is a free function or function group.
\texttt{Poly} is the atom --- everything else is built from it
or depends on it.
\texttt{BezierExpansion} is the molecule ---
it assembles \texttt{Poly} objects into a complete curve.
The print and formatting functions sit on top and never modify
any data, they only read and display.
\end{notebox}

% =============================================================
\needspace{8\baselineskip}
\section{The \texttt{Poly} Struct and Its Constructor}
% =============================================================

\needspace{5\baselineskip}
\subsection{The struct itself}

\begin{minted}{julia}
struct Poly
    coeffs::Vector{Float64}
end
\end{minted}

This defines a new data type called \cd{Poly}.
It has one field: \cd{coeffs}, which must be a
\cd{Vector{Float64}} --- a list of 64-bit floating-point numbers.

To create a \cd{Poly} you would write:

\begin{minted}{julia}
p = Poly([1.0, -2.0, 1.0])
\end{minted}

That works fine. But if you try to pass integers:

\begin{minted}{julia}
p = Poly([1, -2, 1])   # ERROR -- these are Int64, not Float64
\end{minted}

Julia would refuse because the field type is \cd{Vector{Float64}}
and \cd{[1, -2, 1]} is \cd{Vector{Int64}}.
That is why the convenience constructor below exists.

% ---
\subsection{The convenience constructor}

\begin{minted}{julia}
Poly(c::Vector{<:Real}) = Poly(Float64.(c))
\end{minted}

This looks like a function, but it has the same name as the type
(\cd{Poly}), so Julia treats it as an \textbf{outer constructor} ---
a factory function that builds a \cd{Poly}.

\textbf{Breaking it down piece by piece:}

\medskip
\cd{c::Vector\{<:Real\}}

This is a \textbf{type annotation} on the argument \cd{c}.
It says: ``\cd{c} must be a vector, and the elements can be
\emph{any subtype of} \cd{Real}.''

\cd{Real} is Julia's abstract type for all real numbers.
Its subtypes include \cd{Int64}, \cd{Float32}, \cd{Float64},
\cd{Rational}, and many others.
The \cd{<:} symbol means ``is a subtype of''.
So \cd{Vector\{<:Real\}} means ``a vector of any kind of real number.''

\begin{examplebox}
\textbf{What counts as \cd{Vector\{<:Real\}}?}

\medskip
\begin{tabular}{ll}
\toprule
Value & Accepted? \\
\midrule
\cd{[1, 2, 3]}         & Yes --- \cd{Vector\{Int64\}}, and \cd{Int64 <: Real} \\
\cd{[1.0, 2.0, 3.0]}   & Yes --- \cd{Vector\{Float64\}}, and \cd{Float64 <: Real} \\
\cd{[1//2, 3//4]}      & Yes --- \cd{Vector\{Rational\}}, and \cd{Rational <: Real} \\
\cd{["a", "b"]}        & No  --- \cd{String} is not a subtype of \cd{Real} \\
\bottomrule
\end{tabular}
\end{examplebox}

\medskip
\cd{Float64.(c)}

The dot (\cd{.}) before the parentheses is Julia's
\textbf{broadcast} syntax.
It means: ``apply the function \cd{Float64()} to
\emph{every element} of \cd{c}.''

\begin{examplebox}
\cd{Float64.([1, -2, 1])}
$\;\to\;$
\cd{[1.0, -2.0, 1.0]}

Each integer is individually converted to a \cd{Float64}.
\end{examplebox}

\textbf{So the full line reads:}
``If you give me a vector of any real numbers, I will convert
every element to \cd{Float64} and hand that to the real \cd{Poly}
constructor.''

This is called a \textbf{sanity constructor} --- it makes the type
convenient to use without forcing the caller to think about
numeric types.

\begin{notebox}
\textbf{Why not just store integers directly?}

The struct uses \cd{Float64} for all arithmetic that follows
(polynomial evaluation, Horner's method, etc.).
Mixing integer and float arithmetic in Julia can be surprising,
so converting everything to \cd{Float64} at construction time
keeps the rest of the code simple and predictable.
\end{notebox}


\needspace{10\baselineskip}
\subsection*{Summary: \texttt{Poly} Functions}

\begin{notebox}
\textbf{All functions associated with \texttt{Poly}:}

\medskip
\begin{tabular}{lp{3.8cm}p{3.2cm}p{3.5cm}}
\toprule
Function & Signature & Returns & Purpose \\
\midrule
Constructor &
  \cd{Poly(c::Vector\{<:Real\})} &
  \cd{Poly} &
  Convert any numeric vector to \cd{Poly\{Float64\}} \\[4pt]
\cd{*} &
  \cd{(Poly, Real)} or \cd{(Real, Poly)} &
  \cd{Poly} &
  Scale every coefficient by a number \\[4pt]
\cd{+} &
  \cd{(Poly, Poly)} &
  \cd{Poly} &
  Add two polynomials term by term \\[4pt]
\cd{(p)(u)} &
  \cd{(Poly, Real)} &
  \cd{Float64} &
  Evaluate polynomial at $u$ via Horner \\[4pt]
\cd{degree} &
  \cd{(Poly)} &
  \cd{Int} &
  Return \cd{length(coeffs) - 1} \\
\bottomrule
\end{tabular}

\medskip
\textit{In an OOP language these would be methods of a \cd{Poly} class.
In Julia they are standalone functions that happen to take a \cd{Poly}
as their first (or only) argument.
Multiple dispatch handles the rest.}
\end{notebox}


% =============================================================
\needspace{8\baselineskip}
\section{Operator Overriding}
% =============================================================

\needspace{5\baselineskip}
\subsection{What is an operator?}

Operators like \cd{*}, \cd{+}, \cd{-} look special, but in Julia
they are just ordinary \textbf{functions} with short names.
When you write \cd{3 * 5}, Julia is actually calling a function
named \cd{*} with arguments \cd{3} and \cd{5}.

These built-in functions live in a module called \cd{Base},
which is Julia's standard library.
So the full name of the multiplication function is \cd{Base.:*}.

The \cd{.:} syntax: the dot accesses something inside a module
(like \cd{Base.something}), and the colon \cd{:} is needed because
\cd{*} is a special symbol, not a plain identifier.
Without the colon, Julia would try to parse \cd{Base.*} as
invalid syntax.

\needspace{5\baselineskip}
\subsection{Why override?}

Julia already knows how to multiply two numbers: \cd{3 * 5 = 15}.
But it has no idea what \cd{Poly([1.0, 2.0]) * 3} should mean ---
you invented the \cd{Poly} type, so you must teach Julia the rule.

\textbf{Operator overriding} means adding a new method to an
existing operator function, so it knows what to do with your custom type.
You are not replacing the original \cd{*} --- you are
\emph{extending} it with an extra case.

\needspace{5\baselineskip}
\subsection{The two multiplication methods}

\begin{minted}{julia}
# Scalar multiplication
Base.:*(p::Poly, s::Real) = Poly(p.coeffs .* s)
Base.:*(s::Real, p::Poly) = p * s
\end{minted}

\needspace{4\baselineskip}
\subsubsection{First method: \texttt{Poly * number}}

\begin{minted}{julia}
Base.:*(p::Poly, s::Real) = Poly(p.coeffs .* s)
\end{minted}

This teaches Julia what to do when the \textbf{left} argument is
a \cd{Poly} and the \textbf{right} argument is any real number.

\cd{p.coeffs .* s} uses broadcast again: multiply every coefficient
in the vector by the scalar \cd{s}.
The result is wrapped in a new \cd{Poly}.

\begin{examplebox}
\cd{Poly([1.0, -2.0, 1.0]) * 3}

\medskip
\cd{p.coeffs .* s}
$= [1.0, -2.0, 1.0] \;.\!\times\; 3$
$= [3.0, -6.0, 3.0]$

Returns \cd{Poly([3.0, -6.0, 3.0])},
representing $3 - 6u + 3u^2$.
\end{examplebox}

\needspace{4\baselineskip}
\subsubsection{Second method: \texttt{number * Poly}}

\begin{minted}{julia}
Base.:*(s::Real, p::Poly) = p * s
\end{minted}

This teaches Julia what to do when the \textbf{left} argument is
a number and the \textbf{right} argument is a \cd{Poly}.

The body is simply \cd{p * s} --- it flips the order and calls
the first method.
No logic is duplicated; we just say ``this is the same as the
other way around.''

\begin{notebox}
Without this second method, writing \cd{3 * poly} would fail.
Julia does not assume that \cd{*} is commutative ---
it only knows what you explicitly tell it.
For types where order matters (like matrix multiplication),
\cd{A * B} and \cd{B * A} can be completely different operations,
so Julia is right to keep them separate.
\end{notebox}

% ---
\subsection{Multiple dispatch --- how Julia chooses which method to call}

When you write \cd{a * b}, Julia looks at the types of
\textbf{both} \cd{a} and \cd{b} together and finds the method
whose type signature is the best match.
This is called \textbf{multiple dispatch}.

\begin{examplebox}
\textbf{Three calls, three different methods:}

\medskip
\begin{tabular}{lll}
\toprule
Expression & Types & Method chosen \\
\midrule
\cd{3 * 5}        & \cd{(Int64, Int64)}   & built-in numeric \cd{*} \\
\cd{poly * 3}     & \cd{(Poly, Int64)}    & first method above \\
\cd{3 * poly}     & \cd{(Int64, Poly)}    & second method above \\
\bottomrule
\end{tabular}
\end{examplebox}

\textbf{Compare with Python or Java (single dispatch):}
In those languages, \cd{a * b} asks only the type of \cd{a}:
``do \emph{you} know how to multiply yourself by \cd{b}?''
The type of \cd{b} is handled as a secondary concern inside that method.
Julia asks about both arguments simultaneously and dispatches
to the most specific matching method.

This is why the two methods \cd{(Poly, Real)} and \cd{(Real, Poly)}
are genuinely separate in Julia --- they have different type signatures
and Julia treats them as distinct choices.


% =============================================================
\needspace{8\baselineskip}
\section{Polynomial Addition}
% =============================================================

\begin{minted}{julia}
function Base.:+(a::Poly, b::Poly)
    n = max(length(a.coeffs), length(b.coeffs))
    ca = vcat(a.coeffs, zeros(n - length(a.coeffs)))
    cb = vcat(b.coeffs, zeros(n - length(b.coeffs)))
    Poly(ca .+ cb)
end
\end{minted}

This overrides \cd{+} for the case where \textbf{both} arguments
are \cd{Poly} --- another example of multiple dispatch, this time
with a two-argument signature \cd{(Poly, Poly)}.

\textbf{The zero-padding problem.}
Two polynomials can have different degrees, meaning their
\cd{coeffs} vectors have different lengths.
You cannot add two vectors of different lengths element-by-element.

The fix:
\begin{enumerate}
  \item Find the longer length \cd{n}.
  \item Extend the shorter vector with zeros at the end
        (\cd{vcat} concatenates two vectors;
        \cd{zeros(k)} makes a vector of \cd{k} zeros).
  \item Add element-by-element with \cd{.+} (broadcast addition).
\end{enumerate}

\begin{examplebox}
Adding $3 - 6u + 3u^2$ and $1 + 2u$:

\medskip
\cd{a.coeffs = [3.0, -6.0, 3.0]} (length 3) \\
\cd{b.coeffs = [1.0,  2.0]}      (length 2)

\medskip
\cd{n = max(3, 2) = 3}

Pad \cd{b}: \cd{vcat([1.0, 2.0], zeros(1)) = [1.0, 2.0, 0.0]}

Add: \cd{[3.0, -6.0, 3.0] .+ [1.0, 2.0, 0.0] = [4.0, -4.0, 3.0]}

Result: \cd{Poly([4.0, -4.0, 3.0])} $= 4 - 4u + 3u^2$
\end{examplebox}


% =============================================================
\needspace{8\baselineskip}
\section{Polynomial Evaluation via Horner's Method}
% =============================================================

\begin{minted}{julia}
function (p::Poly)(u::Real)
    result = 0.0
    for c in reverse(p.coeffs)
        result = result * u + c
    end
    result
end
\end{minted}

\needspace{5\baselineskip}
\subsection{Callable structs}

The function signature \cd{(p::Poly)(u::Real)} looks unusual.
It means: ``make a \cd{Poly} object callable like a function.''

After this definition, if \cd{p} is a \cd{Poly}, you can write
\cd{p(0.5)} and Julia will call this method, passing \cd{p} as the
struct and \cd{0.5} as \cd{u}.
This is called a \textbf{functor} --- an object that behaves like a function.

\needspace{5\baselineskip}
\subsection{Horner's method}

To evaluate $c_0 + c_1 u + c_2 u^2 + c_3 u^3$, the naive approach
computes each power separately: $u^2$, $u^3$, etc.
Horner's method rewrites the same polynomial as nested multiplications:

\[
  c_0 + u\bigl(c_1 + u(c_2 + u \cdot c_3)\bigr)
\]

This requires only one multiplication and one addition per coefficient,
working from the \emph{highest} degree down to the constant term.
The loop iterates over \cd{reverse(p.coeffs)}, which goes from the
highest-degree coefficient down to the constant.

\begin{examplebox}
Evaluate \cd{Poly([1.0, 4.0, -4.0])} at $u = 0.5$
(this is $1 + 4u - 4u^2$):

\medskip
\cd{reverse([1.0, 4.0, -4.0]) = [-4.0, 4.0, 1.0]}

\medskip
\begin{tabular}{lll}
\toprule
Iteration & \cd{c} & \cd{result = result * u + c} \\
\midrule
start  & ---    & \cd{result = 0.0} \\
1      & \cd{-4.0} & \cd{0.0 * 0.5 + (-4.0) = -4.0} \\
2      & \cd{4.0}  & \cd{-4.0 * 0.5 + 4.0 = 2.0} \\
3      & \cd{1.0}  & \cd{2.0 * 0.5 + 1.0 = 2.0} \\
\bottomrule
\end{tabular}

\medskip
Output: \cd{2.0}

Verify: $1 + 4(0.5) - 4(0.5)^2 = 1 + 2 - 1 = 2.0$ \checkmark
\end{examplebox}




% =============================================================
\needspace{8\baselineskip}
\section{The \texttt{expand\_bernstein} Function and Pascal's Triangle}
% =============================================================

\textit{This is where the real algebra happens.
Everything before this was setup.
Now we cook.}

\begin{minted}{julia}
function expand_bernstein(n::Int, i::Int)::Poly
    m = n - i
    one_minus_u = Poly([binomial(m, j) * (-1)^j for j in 0:m])
    shifted = Poly(vcat(zeros(i), one_minus_u.coeffs))
    shifted * binomial(n, i)
end
\end{minted}

This function takes two integers --- the degree $n$ and the index $i$
--- and returns a \cd{Poly} representing the Bernstein basis polynomial
$B(n, i, u) = \binom{n}{i}\,u^i\,(1-u)^{n-i}$ fully expanded
into monomial form.

The key ingredient is \cd{binomial(n, k)}, Julia's built-in
function for the binomial coefficient $\binom{n}{k}$.
Before we see how \cd{expand\_bernstein} uses it,
let us understand what it actually computes.

% ---
\subsection{What is \texttt{binomial(n, k)}?}

\cd{binomial(n, k)} computes $\binom{n}{k}$,
read aloud as ``$n$ choose $k$.''
It counts the number of ways to choose $k$ items
from a set of $n$ items, without caring about order.

The formula is:
\[
  \binom{n}{k} = \frac{n!}{k!\,(n-k)!}
\]

\begin{examplebox}
\textbf{A few values:}

\medskip
\cd{binomial(4, 0) = 1} \quad (one way to choose nothing) \\
\cd{binomial(4, 1) = 4} \quad (four ways to choose one item from four) \\
\cd{binomial(4, 2) = 6} \quad ($\frac{4!}{2!\,2!} = \frac{24}{4} = 6$) \\
\cd{binomial(4, 3) = 4} \\
\cd{binomial(4, 4) = 1} \quad (one way to choose everything)

\medskip
Notice the symmetry: \cd{binomial(4,1) == binomial(4,3)}.
In general, $\binom{n}{k} = \binom{n}{n-k}$.
\end{examplebox}

% ---
\subsection{Pascal's Triangle --- the Cheat Sheet Nobody Printed}

\textit{Your teacher told you to memorise this.
You didn't. Neither did we.
But now it actually matters.}

Pascal's triangle is a visual arrangement of all binomial coefficients.
Each row $n$ contains $\binom{n}{0}, \binom{n}{1}, \ldots, \binom{n}{n}$.
Each entry is the sum of the two entries directly above it.

\[
\begin{array}{ccccccccc}
n=0: & & & & 1 & & & & \\
n=1: & & & 1 & & 1 & & & \\
n=2: & & 1 & & 2 & & 1 & & \\
n=3: & 1 & & 3 & & 3 & & 1 & \\
n=4: & 1 & & 4 & & 6 & & 4 & 1 \\
\end{array}
\]

For this program, row $n=4$ is the one that matters (quartic case).
Reading left to right: $1, 4, 6, 4, 1$.

These are precisely the prefactors of the five Bernstein
basis polynomials $B(4,0)$ through $B(4,4)$:
\[
  B(4,0) = 1\cdot(1-u)^4, \quad
  B(4,1) = 4\cdot u(1-u)^3, \quad
  B(4,2) = 6\cdot u^2(1-u)^2, \quad \ldots
\]

\begin{notebox}
\textbf{The rule for building the triangle:}
any entry equals the sum of the two entries above it.
For example, $6 = 3 + 3$, and $4 = 1 + 3$.
This means you never need the factorial formula ---
just add neighbours.
Of course, Julia's \cd{binomial()} handles this for you,
so you don't need to add anything by hand either.
\end{notebox}

% ---
\subsection{How \texttt{expand\_bernstein} Uses \texttt{binomial}}

The function expands $B(n,i,u)$ in three steps.
Let us trace through \cd{expand\_bernstein(4, 2)},
which should give us $B(4,2,u) = 6u^2(1-u)^2$.

\textbf{Step A --- expand $(1-u)^m$ using the binomial theorem.}

With $n=4$ and $i=2$, we have $m = n - i = 2$.

The binomial theorem says:
\[
  (1-u)^m = \sum_{j=0}^{m} \binom{m}{j}(-1)^j\,u^j
\]

In code:
\begin{minted}{julia}
one_minus_u = Poly([binomial(m, j) * (-1)^j for j in 0:m])
\end{minted}

This is a \textbf{list comprehension}: Julia evaluates
\cd{binomial(m,j) * (-1)\^{}j} for each \cd{j} from \cd{0} to \cd{m},
and collects the results into a vector.

\begin{examplebox}
For $m = 2$, iterating \cd{j = 0, 1, 2}:

\medskip
\begin{tabular}{llll}
\toprule
\cd{j} & \cd{binomial(2,j)} & \cd{(-1)\^{}j} & product \\
\midrule
0 & 1 & $+1$ & $+1$ \\
1 & 2 & $-1$ & $-2$ \\
2 & 1 & $+1$ & $+1$ \\
\bottomrule
\end{tabular}

\medskip
\cd{one\_minus\_u = Poly([1.0, -2.0, 1.0])} $= 1 - 2u + u^2$

This is $(1-u)^2$. \checkmark
\end{examplebox}

\textbf{Step B --- multiply by $u^i$ by shifting coefficients.}

Multiplying a polynomial by $u^i$ shifts all its powers up by $i$.
In coefficient-vector terms, that means prepending $i$ zeros.

\begin{minted}{julia}
shifted = Poly(vcat(zeros(i), one_minus_u.coeffs))
\end{minted}

\cd{vcat} concatenates two vectors vertically (end to end).

\begin{examplebox}
For $i = 2$: prepend 2 zeros to \cd{[1.0, -2.0, 1.0]}:
\[
  \cd{vcat([0.0, 0.0], [1.0, -2.0, 1.0])}
  = \cd{[0.0, 0.0, 1.0, -2.0, 1.0]}
\]
This represents $u^2 - 2u^3 + u^4$,
which is indeed $(1-u)^2$ multiplied by $u^2$. \checkmark
\end{examplebox}

\textbf{Step C --- multiply by $\binom{n}{i}$.}

\begin{minted}{julia}
shifted * binomial(n, i)
\end{minted}

This uses the scalar \cd{*} operator we defined earlier.

\begin{examplebox}
$\binom{4}{2} = 6$.

\cd{Poly([0.0, 0.0, 1.0, -2.0, 1.0]) * 6}
$= $ \cd{Poly([0.0, 0.0, 6.0, -12.0, 6.0])}

This is $6u^2 - 12u^3 + 6u^4$.

Which matches the output from the program:
\begin{minted}{text}
  B(4,2,u)  =  6u^2(1-u)^2
           =  6u^2 - 12u^3 + 6u^4
\end{minted}
\checkmark
\end{examplebox}

\begin{notebox}
\textbf{Why does the alternating sign $(-1)^j$ appear?}

The binomial theorem for $(a + b)^m$ uses $+$ signs throughout.
But we have $(1 + (-u))^m$, so $b = -u$.
Each power $b^j = (-u)^j = (-1)^j u^j$.
That is where the alternating $+1, -1, +1, \ldots$ pattern comes from.
Without it we would be expanding $(1+u)^m$ instead of $(1-u)^m$,
and the curve would go somewhere completely wrong.

\textit{The sign matters. It always does.}
\end{notebox}

% =============================================================
\needspace{8\baselineskip}
\section{The \texttt{BezierExpansion} Struct --- A Concrete Map}
% =============================================================

Before looking at the print functions, it is important to know
exactly what data lives inside a \cd{BezierExpansion} object.
Every print function receives one as its argument \cd{exp},
so understanding the struct means understanding what every
\cd{exp.field} reference refers to.

\begin{minted}{julia}
struct BezierExpansion
    points ::Matrix{Float64}   # control points (n_pts x dim)
    axes   ::Vector{String}
    n      ::Int               # number of control points
    degree ::Int
    dim    ::Int
    basis  ::Vector{Poly}      # expanded Bernstein poly per point
    polys  ::Dict{String,Poly} # final collected poly per axis
end
\end{minted}

Below is the struct filled in with the actual values for
\textbf{Case 1} (quadratic 2-D, 3 control points).
This is the object that gets passed to every print function.

\medskip
\begin{center}
\begin{tabular}{llp{7.5cm}}
\toprule
Field & Type & Value for Case 1 \\
\midrule
\cd{points} & \cd{Matrix\{Float64\}} &
  $\begin{bmatrix}1 & 1 \\ 3 & 3 \\ 5 & 1\end{bmatrix}$
  Row = one point. Column 1 = $x$, column 2 = $y$. \\[10pt]
\cd{axes} & \cd{Vector\{String\}} &
  \cd{["x", "y"]} \\[4pt]
\cd{n} & \cd{Int} &
  \cd{3} (number of control points) \\[4pt]
\cd{degree} & \cd{Int} &
  \cd{2} ($= n - 1$) \\[4pt]
\cd{dim} & \cd{Int} &
  \cd{2} (number of axes) \\[4pt]
\cd{basis} & \cd{Vector\{Poly\}} &
  \cd{[Poly([1,-2,1]),} $B(2,0)=1-2u+u^2$ \\
  & & \cd{\ Poly([0,2,-2]),} $B(2,1)=2u-2u^2$ \\
  & & \cd{\ Poly([0,0,1])]} $B(2,2)=u^2$ \\[4pt]
\cd{polys} & \cd{Dict\{String,Poly\}} &
  \cd{"x" => Poly([1,4,0])} $= 1+4u$ \\
  & & \cd{"y" => Poly([1,4,-4])} $= 1+4u-4u^2$ \\
\bottomrule
\end{tabular}
\end{center}

\begin{notebox}
\textbf{Two different polynomial collections:}

\cd{basis} contains the \emph{Bernstein} polynomials --- one per
control point, independent of which curve or axis you are looking at.
These are the building blocks.

\cd{polys} contains the \emph{final} polynomial per axis,
after weighting each basis polynomial by its coordinate and summing.
These are what you actually evaluate to get points on the curve.
\end{notebox}


\needspace{10\baselineskip}
\subsection*{Summary: \texttt{BezierExpansion} Functions}

\begin{notebox}
\textbf{All functions associated with \texttt{BezierExpansion}:}

\medskip
\begin{tabular}{lp{4.5cm}p{2.8cm}p{3.5cm}}
\toprule
Function & Signature & Returns & Purpose \\
\midrule
Constructor &
  \cd{BezierExpansion(points, axes)} &
  \cd{BezierExpansion} &
  Build the full expansion: compute all basis polys and collect per-axis polynomials \\[4pt]
\cd{evaluate} &
  \cd{(BezierExpansion, Real)} &
  \cd{Vector\{Float64\}} &
  Evaluate the curve at $u$; returns one coordinate per axis \\
\bottomrule
\end{tabular}

\medskip
\textit{The constructor does all the heavy lifting --- it calls
\cd{expand\_bernstein} for each control point and accumulates the
weighted sums.
\cd{evaluate} is lightweight by design: it just reads the already-built
\cd{polys} dict and calls Horner on each one.}
\end{notebox}


% =============================================================
\needspace{8\baselineskip}
\section{The Print Functions}
% =============================================================

\needspace{10\baselineskip}
\subsection*{Summary: Formatting Helper Functions}

\begin{notebox}
\textbf{Standalone formatting utilities (not tied to any struct):}

\medskip
\begin{tabular}{lp{4cm}p{2cm}p{4cm}}
\toprule
Function & Signature & Returns & Purpose \\
\midrule
\cd{ruler} &
  \cd{(Char)} &
  \cd{String} &
  64-character divider line \\[4pt]
\cd{fmt\_coeff} &
  \cd{(Float64)} &
  \cd{String} &
  Integer if whole, 4 d.p.\ otherwise \\[4pt]
\cd{fmt\_poly} &
  \cd{(Poly, String)} &
  \cd{String} &
  Readable polynomial string, e.g.\ \cd{"1 + 4u - 4u\^{}2"} \\[4pt]
\cd{fmt\_bernstein\_symbolic} &
  \cd{(Int, Int)} &
  \cd{String} &
  Symbolic basis string, e.g.\ \cd{"6u\^{}2(1-u)\^{}2"} \\
\bottomrule
\end{tabular}

\medskip
\textit{These functions only produce strings --- they never touch a
\cd{Poly} struct's internals directly, and they never print anything.
Separating formatting from printing means you can reuse them anywhere.}
\end{notebox}

\needspace{5\baselineskip}
\subsection*{The Five Step Functions}

All five print functions receive the same \cd{exp} struct
described above, plus sometimes a list of $u$ values.
None of them return a value --- they only write to the terminal.

% ---
\subsection{\texttt{print\_step1\_bernstein\_form(exp)}}

\textbf{What it prints:} the curve formula in Bernstein form ---
the symbolic equation $\mathbf{P}(u) = B(n,0)P_0 + \cdots$
and the definition of each basis function.

\begin{minted}{julia}
function print_step1_bernstein_form(exp::BezierExpansion)
    print_step_header(1, "Bernstein Form")
    n   = exp.degree
    dim = join(["$(ax)(u)" for ax in exp.axes], "  ")
    println("  P(u) =  [$dim]")
    println()
    terms = ["B($n,$i)*P$i" for i in 0:exp.n-1]
    println("  P(u) =  " * join(terms, "  +  "))
    println()
    println("  where:")
    for i in 0:exp.n-1
        sym = fmt_bernstein_symbolic(n, i)
        println("    B($n,$i,u)  =  $sym")
    end
end
\end{minted}

\textbf{How it works:}

\begin{enumerate}
  \item Reads \cd{exp.degree} (e.g.\ \cd{2}) and \cd{exp.axes}
        (e.g.\ \cd{["x","y"]}).
  \item Builds the axis label string \cd{"x(u)  y(u)"}
        by looping over \cd{exp.axes} and joining with spaces.
  \item Builds the sum string \cd{"B(2,0)*P0  +  B(2,1)*P1  +  B(2,2)*P2"}
        by looping \cd{i} from \cd{0} to \cd{exp.n-1}.
  \item Loops again over the same range and calls
        \cd{fmt\_bernstein\_symbolic(n, i)} to produce the
        human-readable form like \cd{"2u(1-u)"} for each basis function.
\end{enumerate}

\begin{examplebox}
Fields used: \cd{exp.degree}, \cd{exp.axes}, \cd{exp.n}

No actual polynomial data is used here --- this step is purely
symbolic, showing the formula \emph{before} any numbers are substituted.
\end{examplebox}

\textbf{Actual output (Case 2):}
\begin{minted}{text}
  STEP 1 - Bernstein Form
  --------------------------------------------------------------
  P(u) =  [x(u)  y(u)  z(u)]

  P(u) =  B(4,0)*P0  +  B(4,1)*P1  +  B(4,2)*P2
                     +  B(4,3)*P3  +  B(4,4)*P4

  where:
    B(4,0,u)  =  (1-u)^4
    B(4,1,u)  =  4u(1-u)^3
    B(4,2,u)  =  6u^2(1-u)^2
    B(4,3,u)  =  4u^3(1-u)
    B(4,4,u)  =  u^4
\end{minted}

% ---
\subsection{\texttt{print\_step2\_substitution(exp)}}

\textbf{What it prints:} the actual coordinate values substituted
into the Bernstein formula --- one line per axis.

\begin{minted}{julia}
function print_step2_substitution(exp::BezierExpansion)
    print_step_header(2, "Substitution of Control Point Coordinates")
    n = exp.degree
    for (dim_idx, axis) in enumerate(exp.axes)
        parts = String[]
        for i in 1:exp.n
            coord = exp.points[i, dim_idx]
            sym   = fmt_bernstein_symbolic(n, i-1)
            c     = fmt_coeff(coord)
            push!(parts, "  $(sym)*$(c)")
        end
        println("  $(axis)(u) =" * join(parts, "  +  "))
    end
end
\end{minted}

\textbf{How it works:}

The function has a \textbf{double loop}: the outer loop iterates
over axes (\cd{"x"}, \cd{"y"}), and the inner loop iterates over
control points (\cd{i = 1} to \cd{exp.n}).

For each combination of axis and control point, it reads
\cd{exp.points[i, dim\_idx]} --- the coordinate of point $i$
on the current axis.

\begin{examplebox}
\textbf{Reading the points matrix for Case 1:}

\medskip
\cd{exp.points =}
$\begin{bmatrix}1 & 1 \\ 3 & 3 \\ 5 & 1\end{bmatrix}$

\medskip
\begin{tabular}{lllll}
\toprule
\cd{i} & \cd{dim\_idx} & \cd{axis} & \cd{exp.points[i, dim\_idx]} & Meaning \\
\midrule
1 & 1 & \cd{"x"} & \cd{1.0} & $x$-coord of $P_0$ \\
2 & 1 & \cd{"x"} & \cd{3.0} & $x$-coord of $P_1$ \\
3 & 1 & \cd{"x"} & \cd{5.0} & $x$-coord of $P_2$ \\
1 & 2 & \cd{"y"} & \cd{1.0} & $y$-coord of $P_0$ \\
2 & 2 & \cd{"y"} & \cd{3.0} & $y$-coord of $P_1$ \\
3 & 2 & \cd{"y"} & \cd{1.0} & $y$-coord of $P_2$ \\
\bottomrule
\end{tabular}

\medskip
Each coordinate is paired with its symbolic basis term
(e.g.\ \cd{"(1-u)\^{}2*1"}) and all terms for one axis
are joined into one line.
\end{examplebox}

Fields used: \cd{exp.degree}, \cd{exp.axes}, \cd{exp.n}, \cd{exp.points}.

\textbf{Actual output (Case 2):}
\begin{minted}{text}
  STEP 2 - Substitution of Control Point Coordinates
  --------------------------------------------------------------
  x(u) =  (1-u)^4*0  +  4u(1-u)^3*0  +  6u^2(1-u)^2*4
                     +  4u^3(1-u)*4   +  u^4*5
  y(u) =  (1-u)^4*0  +  4u(1-u)^3*4  +  6u^2(1-u)^2*0
                     +  4u^3(1-u)*4   +  u^4*4
  z(u) =  (1-u)^4*1  +  4u(1-u)^3*1  +  6u^2(1-u)^2*1
                     +  4u^3(1-u)*1   +  u^4*1

  (Each basis term will now be expanded individually)
\end{minted}

Notice that the first two $x$ terms are multiplied by 0
(because $P_0$ and $P_1$ have $x=0$), so they will vanish
in Step~4.
All $z$ terms are multiplied by 1 (every point has $z=1$),
which is why $z(u)$ will collapse to the constant~1.

% ---
\subsection{\texttt{print\_step3\_expanded\_terms(exp)}}

\textbf{What it prints:} each Bernstein basis polynomial expanded
into monomial form.

\begin{minted}{julia}
function print_step3_expanded_terms(exp::BezierExpansion)
    print_step_header(3, "Expand Each Bernstein Basis Term")
    n = exp.degree
    for i in 0:exp.n-1
        basis = exp.basis[i+1]
        sym   = fmt_bernstein_symbolic(n, i)
        ex    = fmt_poly(basis)
        println("  B($n,$i,u)  =  $sym")
        println("           =  $ex")
        println()
    end
end
\end{minted}

\textbf{How it works:}

The loop runs \cd{i} from \cd{0} to \cd{exp.n-1} (0-indexed, matching
the mathematical notation $B(n, 0)$ through $B(n, n)$).

Inside the loop:
\begin{itemize}
  \item \cd{exp.basis[i+1]} fetches the $i$-th expanded basis polynomial.
        The \cd{+1} is needed because Julia arrays are 1-indexed
        (the first element is at index 1, not 0).
  \item \cd{fmt\_bernstein\_symbolic(n, i)} produces the symbolic
        string like \cd{"2u(1-u)"}.
  \item \cd{fmt\_poly(basis)} converts the \cd{Poly} coefficients
        into a readable string like \cd{"2u - 2u\^{}2"}.
\end{itemize}

\begin{examplebox}
For Case 1, \cd{exp.basis} contains three \cd{Poly} objects:

\medskip
\begin{tabular}{lll}
\toprule
\cd{i} & \cd{exp.basis[i+1]} & Printed as \\
\midrule
0 & \cd{Poly([1.0, -2.0, 1.0])} & \cd{1 - 2u + u\^{}2} \\
1 & \cd{Poly([0.0,  2.0, -2.0])} & \cd{2u - 2u\^{}2} \\
2 & \cd{Poly([0.0,  0.0,  1.0])} & \cd{u\^{}2} \\
\bottomrule
\end{tabular}
\end{examplebox}

Fields used: \cd{exp.degree}, \cd{exp.n}, \cd{exp.basis}.

\textbf{Actual output (Case 2):}
\begin{minted}{text}
  STEP 3 - Expand Each Bernstein Basis Term
  --------------------------------------------------------------
  B(4,0,u)  =  (1-u)^4
           =  1 - 4u + 6u^2 - 4u^3 + u^4

  B(4,1,u)  =  4u(1-u)^3
           =  4u - 12u^2 + 12u^3 - 4u^4

  B(4,2,u)  =  6u^2(1-u)^2
           =  6u^2 - 12u^3 + 6u^4

  B(4,3,u)  =  4u^3(1-u)
           =  4u^3 - 4u^4

  B(4,4,u)  =  u^4
           =  u^4
\end{minted}

Each pair of lines is one iteration of the loop:
the first line comes from \cd{fmt\_bernstein\_symbolic},
the second from \cd{fmt\_poly(exp.basis[i+1])}.
The \cd{exp.basis} vector holds these five \cd{Poly} objects
(one per control point), already computed by \cd{expand\_bernstein}
during construction.

% ---
\subsection{\texttt{print\_step4\_collected(exp)}}

\textbf{What it prints:} for each axis, shows each weighted
contribution (\cd{basis[i] * coord}), then the final
collected polynomial.

\begin{minted}{julia}
function print_step4_collected(exp::BezierExpansion)
    print_step_header(4, "Collect Terms - Final Monomial Polynomial")
    n = exp.degree
    for (dim_idx, axis) in enumerate(exp.axes)
        println("  $(axis)(u) - weighted contributions:")
        for i in 1:exp.n
            coord   = exp.points[i, dim_idx]
            term    = exp.basis[i] * coord
            exp_str = fmt_poly(term)
            c_str   = fmt_coeff(coord)
            println("    B($n,$(i-1))*$(lpad(c_str,4))  =  $exp_str")
        end
        println()
    end
    println("  Collected:")
    for axis in exp.axes
        poly = exp.polys[axis]
        println("    $(axis)(u)  =  $(fmt_poly(poly))")
    end
end
\end{minted}

\textbf{How it works:}

Again a double loop: outer over axes, inner over control points.

Inside the inner loop:
\begin{itemize}
  \item \cd{exp.points[i, dim\_idx]} reads the coordinate (same as Step 2).
  \item \cd{exp.basis[i] * coord} multiplies the expanded basis
        polynomial by that coordinate using the \cd{*} override
        we defined earlier.
        This produces a new \cd{Poly} --- the weighted contribution.
  \item \cd{fmt\_poly(term)} converts that \cd{Poly} to a string
        for printing.
  \item \cd{lpad(c\_str, 4)} left-pads the coordinate string to
        4 characters so the columns line up neatly.
\end{itemize}

After both loops, a second loop over \cd{exp.axes} prints the
\emph{final} collected polynomial from \cd{exp.polys[axis]}.
This is the result that was computed once during construction
and stored in the struct --- the print function just reads it.

\begin{examplebox}
\textbf{For the $x$ axis in Case 1:}

\medskip
\begin{tabular}{llll}
\toprule
\cd{i} & coord & \cd{basis[i] * coord} & Printed \\
\midrule
1 & 1.0 & \cd{Poly([1,-2,1])*1.0} & \cd{1 - 2u + u\^{}2} \\
2 & 3.0 & \cd{Poly([0,2,-2])*3.0} & \cd{6u - 6u\^{}2} \\
3 & 5.0 & \cd{Poly([0,0,1])*5.0}  & \cd{5u\^{}2} \\
\bottomrule
\end{tabular}

\medskip
Then: \cd{exp.polys["x"] = Poly([1.0, 4.0, 0.0])} is printed as \cd{1 + 4u}.
\end{examplebox}

Fields used: \cd{exp.degree}, \cd{exp.axes}, \cd{exp.n},
\cd{exp.points}, \cd{exp.basis}, \cd{exp.polys}.

\textbf{Actual output (Case 2):}
\begin{minted}{text}
  STEP 4 - Collect Terms - Final Monomial Polynomial
  --------------------------------------------------------------
  x(u) - weighted contributions:
    B(4,0)*   0  =  0
    B(4,1)*   0  =  0
    B(4,2)*   4  =  24u^2 - 48u^3 + 24u^4
    B(4,3)*   4  =  16u^3 - 16u^4
    B(4,4)*   5  =  5u^4

  y(u) - weighted contributions:
    B(4,0)*   0  =  0
    B(4,1)*   4  =  16u - 48u^2 + 48u^3 - 16u^4
    B(4,2)*   0  =  0
    B(4,3)*   4  =  16u^3 - 16u^4
    B(4,4)*   4  =  4u^4

  z(u) - weighted contributions:
    B(4,0)*   1  =  1 - 4u + 6u^2 - 4u^3 + u^4
    B(4,1)*   1  =  4u - 12u^2 + 12u^3 - 4u^4
    B(4,2)*   1  =  6u^2 - 12u^3 + 6u^4
    B(4,3)*   1  =  4u^3 - 4u^4
    B(4,4)*   1  =  u^4

  Collected:
    x(u)  =  24u^2 - 32u^3 + 13u^4
    y(u)  =  16u - 48u^2 + 64u^3 - 28u^4
    z(u)  =  1
\end{minted}

The inner loop produces one row per control point per axis.
The \cd{lpad(c\_str, 4)} call right-aligns the coordinate
(e.g.\ \cd{"   0"}, \cd{"   4"}, \cd{"   5"}) so the equals
signs line up.
After the inner loop ends, \cd{exp.polys[axis]} is read directly
for the ``Collected'' lines --- that value was computed once in
the constructor and is just retrieved here.

% ---
\subsection{\texttt{print\_step5\_evaluation(exp, u\_values)}}

\textbf{What it prints:} a table of curve points at each $u$ value,
then a detailed substitution check at the middle $u$.

\begin{minted}{julia}
function print_step5_evaluation(exp::BezierExpansion,
                                u_values::Vector{Float64})
    print_step_header(5, "Evaluation at Given u Values")
    axes  = exp.axes
    dim   = length(axes)
    col_u = 6
    col_v = 10

    header = "  " * lpad("u", col_u)
    for ax in axes
        header *= "  " * lpad("$(ax)(u)", col_v)
    end
    println(header)
    println("  " * "-"^(col_u + (col_v + 2) * dim + 2))

    for u in u_values
        pt  = evaluate(exp, u)
        row = "  " * lpad(@sprintf("%.2f", u), col_u)
        for v in pt
            row *= "  " * lpad(@sprintf("%.6f", v), col_v)
        end
        println(row)
    end
    ...
end
\end{minted}

\textbf{How it works:}

\textbf{The header row.}
\cd{lpad("u", 6)} pads the string \cd{"u"} to width 6
by adding spaces on the left.
The same is done for each axis label.
This creates aligned columns.

\textbf{The data rows.}
For each \cd{u} in \cd{u\_values}, it calls \cd{evaluate(exp, u)},
which returns a vector with one value per axis
(e.g.\ \cd{[3.0, 2.0]} for $x$ and $y$).
Each value is formatted to 6 decimal places with \cd{@sprintf}
and padded to width 10 with \cd{lpad}.

\textbf{The verification block.}
After the table, the code picks the middle \cd{u}:

\begin{minted}{julia}
u_check = u_values[div(length(u_values), 2) + 1]
\end{minted}

\cd{div} is integer division (no remainder).
For 5 values: \cd{div(5, 2) = 2}, so \cd{u\_values[3] = 0.50}.

Then for each axis it loops over the polynomial's coefficients
and prints each term as a separate substitution
(e.g.\ \cd{"4*0.50\^{}1"}), so you can follow the arithmetic by hand.

\begin{examplebox}
\textbf{Fields and arguments used:}

\medskip
\begin{tabular}{ll}
\toprule
What & Where it comes from \\
\midrule
Axis labels in header & \cd{exp.axes} \\
Curve values in rows  & \cd{evaluate(exp, u)} $\to$ \cd{exp.polys} \\
Middle $u$ pick       & \cd{u\_values} (the external argument) \\
Substitution terms    & \cd{exp.polys[axis].coeffs} looped directly \\
\bottomrule
\end{tabular}
\end{examplebox}

\textbf{Actual output (Case 2):}
\begin{minted}{text}
  STEP 5 - Evaluation at Given u Values
  --------------------------------------------------------------
       u        x(u)        y(u)        z(u)
  --------------------------------------------
    0.00    0.000000    0.000000    1.000000
    0.25    1.050781    1.890625    1.000000
    0.50    2.812500    2.250000    1.000000
    0.75    4.113281    3.140625    1.000000
    1.00    5.000000    4.000000    1.000000

  Verification at u = 0.50:
    x(0.50)  =  24u^2 - 32u^3 + 13u^4
               =  24*0.50^2  +  -32*0.50^3  +  13*0.50^4
               =  2.812500

    y(0.50)  =  16u - 48u^2 + 64u^3 - 28u^4
               =  16*0.50^1  +  -48*0.50^2  +  64*0.50^3  +  -28*0.50^4
               =  2.250000

    z(0.50)  =  1
               =  1
               =  1.000000
\end{minted}

The table is built column by column: \cd{u} is printed first
(width 6), then each axis value (width 10), all left-padded
so decimal points align.
The verification block loops over \cd{poly.coeffs} directly,
skipping any coefficient whose absolute value is below $10^{-9}$
(same threshold used in \cd{fmt\_poly}).
For $z(u) = 1$ there is only a constant term, so the
substitution line just shows \cd{"= 1"} with no powers of $u$.



\needspace{12\baselineskip}
\subsection*{Summary: Print Functions}

\begin{notebox}
\textbf{All print functions (none return a value --- output only):}

\medskip
\begin{tabular}{lp{4.8cm}p{5cm}}
\toprule
Function & Signature & Purpose \\
\midrule
\cd{print\_step\_header} &
  \cd{(Int, String)} &
  Print numbered step heading with dashed rule \\[4pt]
\cd{print\_title} &
  \cd{(BezierExpansion, Int)} &
  Print the title banner with degree and dimension info \\[4pt]
\cd{print\_control\_points} &
  \cd{(BezierExpansion)} &
  List all control points with coordinates \\[4pt]
\cd{print\_step1\_bernstein\_form} &
  \cd{(BezierExpansion)} &
  Show symbolic Bernstein formula and basis definitions \\[4pt]
\cd{print\_step2\_substitution} &
  \cd{(BezierExpansion)} &
  Show coordinates substituted into basis terms \\[4pt]
\cd{print\_step3\_expanded\_terms} &
  \cd{(BezierExpansion)} &
  Show each basis polynomial in monomial form \\[4pt]
\cd{print\_step4\_collected} &
  \cd{(BezierExpansion)} &
  Show weighted contributions and final collected polynomial \\[4pt]
\cd{print\_step5\_evaluation} &
  \cd{(BezierExpansion, Vector\{Float64\})} &
  Print evaluation table and mid-point verification \\[4pt]
\cd{print\_all} &
  \cd{(BezierExpansion, Vector\{Float64\}, Int)} &
  Orchestrate steps 1--5 in order \\
\bottomrule
\end{tabular}

\medskip
\textit{\cd{print\_all} is the conductor.
It calls everything in sequence and adds nothing of its own.
If you ever want to print just one step in isolation ---
say, only the evaluation table --- you can call \cd{print\_step5}
directly. The steps are independent.}
\end{notebox}


% =============================================================
\needspace{8\baselineskip}
\section{Summary: How Data Flows Through the Program}
% =============================================================

\begin{center}
\begin{tabular}{rp{4cm}p{6cm}}
\toprule
Step & Function & Data used from \cd{exp} \\
\midrule
1 & \cd{print\_step1} &
  \cd{degree}, \cd{axes}, \cd{n} --- symbolic only, no numbers \\
2 & \cd{print\_step2} &
  \cd{degree}, \cd{axes}, \cd{n}, \cd{points} --- coordinates substituted \\
3 & \cd{print\_step3} &
  \cd{degree}, \cd{n}, \cd{basis} --- expanded basis polys printed \\
4 & \cd{print\_step4} &
  \cd{degree}, \cd{axes}, \cd{n}, \cd{points}, \cd{basis}, \cd{polys} ---
  weighted terms and final result \\
5 & \cd{print\_step5} &
  \cd{axes}, \cd{polys} --- final polys evaluated at each $u$ \\
\bottomrule
\end{tabular}
\end{center}

\vspace{2em}\hrule\vspace{0.5em}
{\small\textcolor{gray}{Code explanation for \texttt{01-bezier-algebra.jl} --- output examples use Case~2 \textbar{} Prepared by E.R.\ Nurwijayadi \textbar{} Prompt-driven, AI-composed}}

\end{document}
